Antes de interpretar los grafos Mapper, fue necesario reconciliar las salidas existentes del repositorio. El problema principal era que los artefactos individuales guardados en \texttt{outputs/mapper/} no coincidían completamente con las tablas agregadas activas usadas por el reporte. En particular, existían archivos para 21 grafos Mapper, pero las tablas agregadas principales resumían únicamente siete grafos con lente \texttt{pca2}. Es decir, había más resultados generados que resultados integrados en las tablas de interpretación.

Esta diferencia es importante porque el reporte no debe basarse en una lectura ambigua de los archivos disponibles. Si un grafo existe como artefacto individual, pero no aparece en las tablas agregadas, no queda claro si forma parte de la interpretación activa o si es un resultado anterior no sincronizado. Por eso se implementó una etapa de reconciliación para distinguir entre artefactos existentes, tablas activas y resultados efectivamente usados para la interpretación científica.

La reconciliación encontró 21 archivos de grafos en formato JSON, 21 archivos de nodos en formato CSV, 21 archivos de aristas en formato CSV y 21 archivos de configuración en formato JSON. Además, se verificó que cada grafo tuviera su archivo correspondiente de nodos, aristas y configuración. Por lo tanto, los artefactos existentes forman 21 conjuntos completos:

\[
\text{graph} + \text{nodes} + \text{edges} + \text{config}.
\]

Estos 21 conjuntos corresponden a la combinación de siete espacios de variables y tres lentes:

\[
7 \text{ espacios de variables} \times 3 \text{ lentes} = 21 \text{ grafos Mapper}.
\]

\begin{table}[H]
\centering
\small
\begin{tabular}{lr}
\toprule
\textbf{Tipo de artefacto} & \textbf{Cantidad encontrada} \\
\midrule
Archivos de grafos JSON & 21 \\
Archivos de nodos CSV & 21 \\
Archivos de aristas CSV & 21 \\
Archivos de configuración JSON & 21 \\
Conjuntos completos graph--nodes--edges--config & 21 \\
Espacios de variables & 7 \\
Lentes & 3 \\
\bottomrule
\end{tabular}
\caption{Resumen de artefactos Mapper encontrados durante la reconciliación de salidas.}
\label{tab:output_reconciliation_counts}
\end{table}

La diferencia central está entre los artefactos existentes y las tablas agregadas activas. Los artefactos existentes son los archivos individuales producidos por las corridas Mapper, ubicados en carpetas como \texttt{outputs/mapper/graphs/}, \texttt{outputs/mapper/nodes/}, \texttt{outputs/mapper/edges/} y \texttt{outputs/mapper/config/}. En cambio, las tablas agregadas activas son resúmenes ya procesados, como métricas de grafos, comparación entre espacios, sensibilidad por lente y selección de grafos principales. La reconciliación mostró que estas tablas activas no estaban sincronizadas con todos los artefactos existentes.

En particular, las tablas agregadas originales estaban restringidas a la capa \texttt{pca2}. Esto sugiere que una etapa posterior del pipeline generó o sobrescribió las tablas principales enfocándose únicamente en los grafos \texttt{pca2}, que son los usados para la interpretación central. Sin embargo, los artefactos individuales conservaban también grafos construidos con las lentes \texttt{domain} y \texttt{density}. Por esta razón, el reporte distingue entre la capa activa de interpretación y las capas disponibles para sensibilidad metodológica.

Para resolver esta desalineación sin borrar ni modificar las tablas originales, se reconstruyeron nuevas tablas con el sufijo \texttt{\_all\_existing}. Estas tablas resumen todos los artefactos encontrados, no solamente los grafos \texttt{pca2}. Los archivos reconstruidos fueron:

\begin{itemize}
    \item \texttt{output\_manifest.csv}
    \item \texttt{output\_consistency\_warnings.md}
    \item \texttt{mapper\_graph\_metrics\_all\_existing.csv}
    \item \texttt{mapper\_lens\_sensitivity\_all\_existing.csv}
    \item \texttt{mapper\_space\_comparison\_all\_existing.csv}
\end{itemize}

El archivo \texttt{output\_manifest.csv} funciona como inventario de los artefactos existentes. Registra qué grafos tienen archivos correspondientes de nodos, aristas y configuración, y permite identificar qué artefactos aparecen o no en las tablas agregadas activas. Por su parte, \texttt{mapper\_graph\_metrics\_all\_existing.csv} reconstruye métricas topológicas para los 21 grafos disponibles, mientras que las tablas de sensibilidad por lente y comparación por espacio permiten revisar de forma más completa las salidas ya generadas.

Un resultado importante de esta reconciliación es que no hay evidencia de una grilla completa de hiperparámetros en \texttt{n\_cubes} y \texttt{overlap}. Aunque el código puede soportar múltiples combinaciones, los artefactos existentes corresponden a una sola configuración de cubierta:

\[
\texttt{n\_cubes}=10,\qquad \texttt{overlap}=0.35.
\]

En los nombres de los grafos esta configuración aparece como \texttt{cubes10\_overlap0p35}. Por lo tanto, el análisis no debe describirse como una búsqueda completa de hiperparámetros. La formulación correcta es que se compararon espacios de variables y lentes bajo una cubierta fija.

\begin{quote}
\textbf{Nota de alcance.} Este reporte no interpreta los resultados como una grilla completa de \texttt{n\_cubes} y \texttt{overlap}. Los artefactos reconciliados permiten comparar siete espacios de variables y tres lentes, pero todos bajo la misma cubierta \texttt{cubes10\_overlap0p35}.
\end{quote}

A partir de esta reconciliación, la capa activa de interpretación del reporte se define como el conjunto de grafos seleccionados con lente \texttt{pca2}, listados en \texttt{outputs/mapper/tables/main\_graph\_selection.csv}. Esta decisión no implica que los grafos \texttt{domain} y \texttt{density} sean irrelevantes. Más bien, estos se tratan como capas de sensibilidad metodológica: sirven para revisar si la estructura observada depende fuertemente del lente usado, pero no constituyen la base principal de las conclusiones.

En resumen, la reconciliación de salidas cumple tres funciones. Primero, confirma que existen 21 grafos Mapper completos. Segundo, aclara que las tablas activas originales estaban enfocadas en \texttt{pca2}. Tercero, delimita el alcance metodológico del reporte: la interpretación científica principal se basa en los grafos \texttt{pca2} seleccionados, mientras que los demás lentes se usan como evidencia complementaria de sensibilidad. Esta separación permite que las conclusiones sean reproducibles y evita sobreafirmar una búsqueda de hiperparámetros que no está materializada en los artefactos actuales.